\scanpage{76}
\mspage{75}
\begin{center}- 33 -\end{center}
\begin{flushright}n°357\end{flushright}
\[
\varepsilon_x*f-\varepsilon_{x_0}*f_0
=\varepsilon_x*(f-f_1)-\varepsilon_{x_0}*(f_0-f_1)+(\varepsilon_x-\varepsilon_{x_0})*f_1
\]
et par suite, si $N_p(f-f_0)\leq\delta$ et $x\in\unsure{K|\mathcal K}$ :
\[
N_p(\varepsilon_x*f-\varepsilon_{x_0}*f_0)\leq3C\delta+N_p((\varepsilon_x-\varepsilon_{x_0})*f_1)
\]
où on a posé $C=\sup_K(\chi^{-1/q})$. Mais on a :
\[
N_p((\varepsilon_x-\varepsilon_{x_0})*f_1)^p
=\int\left|f_1(x^{-1}z)\chi(x^{-1})-f_1(x_0^{-1}z)\chi(x_0^{-1})\right|^p\,d\gamma(z)\ ;
\]
si $K_1$ est le support de $f_1$, la fonction $f_1(x^{-1}z)\chi(x^{-1})$ induit sur $K\times E$ une fonction continue dont le support est contenu dans $K\times(K.K_1)$ ; celle-ci est donc uniformément continue ; on en conclut aussitôt qu'il existe un voisinage $V$ de $x_0$ dans $K$ tel que le second membre de la relation ci-dessus soit $\leq\delta$ pour $x\in V$. Cela achève la démonstration.

\uline{Remarque} 1.- Le formule (10) reste évidemment valable pour $p=+\infty$ ; la démonstration ci-dessus montre alors que $(x,f)\to\varepsilon_x*f$ est une application continue de $G\times\overline{\mathcal K(E)}$ dans $\overline{\mathcal K(E)}$.

\uline{Remarque} 2.- Comme $\varepsilon_x*f$ est négligeable si $f$ l'est [N.B. - Pour le style, cf. ``puisque mon époux l'est'', P. Corneille, \uline{Horace}, \unsure{1e|le} éd., Acte I, Sc. 1], l'application $(x,f)\to\varepsilon_x*f$ détermine, par passage au quotient, une application de $G\times L^p$ dans $L^p$, qu'on note aussi $(x,f)\to\varepsilon_x*f$, et qui satisfait aussi à (10). Une remarque analogue s'applique aux opérations définies ci-dessous.

PROPOSITION 8.- \uline{Soit} $\lambda$ \uline{une mesure bornée sur} $G$ ; \uline{pour} $1\leq p<+\infty$, \uline{soit} $\mu=\chi^{1/q}.\lambda$, \uline{et soit} $f\in\mathcal L^p(E,\gamma)$. \uline{Alors} $\mu$ \uline{et} $f$ \uline{sont composables, et on a}
\[
\tag{11}
N_p((\chi^{1/q}.\lambda)*f)\leq N_p(f).\lVert\lambda\rVert.
\]
Le lemme 2 du n°2 montre que le support de $\lambda$, donc celui de $\mu$, est réunion dénombrable de compacts ; on peut donc appliquer la prop.6 du n°2, où l'on a, d'après la prop.7 et l'inégalité de Hölder :
